\begin{table*}[t]
\centering
\caption{固定网络权重的步长加密结果。杆端误差单位为 $\mu\mathrm{m}$，报告五种子均值与样本标准差。}
\label{tab:neural_refinement}
\small
\begin{tabular}{llrrrr}
\toprule
比较对象 & 模型 & $h$ & $h/2$ & $h/4$ & $h/8$ \\
\midrule
连续参考 & $H_0$（无网络） & 77.466 $\pm$ 8.552 & 18.836 $\pm$ 2.048 & 4.350 $\pm$ 0.442 & 1.563 $\pm$ 0.342 \\
连续参考 & GR & 78.263 $\pm$ 8.899 & 19.596 $\pm$ 2.396 & 4.891 $\pm$ 0.777 & 1.335 $\pm$ 0.342 \\
连续参考 & GC & 78.298 $\pm$ 8.670 & 19.628 $\pm$ 2.157 & 4.908 $\pm$ 0.522 & 1.287 $\pm$ 0.129 \\
\midrule
同步长教师 & $H_0$（无网络） & 1.645 $\pm$ 0.471 & 1.647 $\pm$ 0.471 & 1.647 $\pm$ 0.471 & 1.647 $\pm$ 0.471 \\
同步长教师 & GR & 0.585 $\pm$ 0.239 & 0.586 $\pm$ 0.239 & 0.586 $\pm$ 0.240 & 0.586 $\pm$ 0.240 \\
同步长教师 & GC & 0.408 $\pm$ 0.225 & 0.408 $\pm$ 0.225 & 0.408 $\pm$ 0.226 & 0.408 $\pm$ 0.226 \\
\bottomrule
\end{tabular}
\par\smallskip\footnotesize 所有模型使用相同的八条测试轨迹、914个输出时刻和实际时长15.993 s；各轨迹先计算时间RMSE，再在种子内平均。$h=0.009999813692$ 为无量纲步长。同步长教师由完整非线性Hamiltonian在相同步长下求解；连续参考保持原DOP853解。最细步长的完整状态相对$L^2$误差为：$H_0$ $0.001353\pm0.000101$，GR $0.002075\pm0.000663$，GC $0.002104\pm0.000202$；该指标中$H_0$更低。R/C的完整结果随复现数据提供。
\end{table*}
